% =====================================================================
\chapter{Rectas, planos, sistemas lineales y determinantes}
\label{tema:2}
% =====================================================================

Una ecuación lineal en tres incógnitas describe un plano, y resolver un sistema de tres ecuaciones es cortar tres planos. Este tema escribe rectas y planos con vectores, resuelve sistemas pequeños por eliminación gaussiana e interpreta el determinante como un área o un volumen con signo. Todo se limita a tamaños $2\times2$ y $3\times3$: la teoría general (rango, independencia lineal, determinantes de orden $n$) es materia de los módulos~2 y~3 del curso. La presentación sigue las sesiones~5--16 de la unidad~1 de 18.02SC~\cite{mit1802}, la §2.5 de OpenStax~\cite{openstax3}, las secciones~11.2--11.4 de Strang~\cite{strangcalc} y, para la eliminación, el capítulo~2 de su \emph{Introduction to Linear Algebra}~\cite{strangla}.

\section{Rectas y planos}

\begin{definicion}[Recta]\label{def:recta}
La \emph{recta} que pasa por el punto $P_0$ con vector director $\vv\neq\vcero$ es el conjunto de puntos
\[
\vr(t)=P_0+t\,\vv,\qquad t\in\R .
\]
En coordenadas, $x=x_0+tv_1$, $y=y_0+tv_2$, $z=z_0+tv_3$ (\emph{ecuaciones paramétricas}). La recta que pasa por $P$ y $Q$ tiene vector director $\overrightarrow{PQ}$.
\end{definicion}

En el plano $\R^2$ una recta admite además una ecuación \emph{implícita} $ax+by=c$; el vector $(a,b)$ es perpendicular a la recta, porque si $(x_1,y_1)$ y $(x_2,y_2)$ son dos de sus puntos, restando las ecuaciones queda $(a,b)\cdot(x_2-x_1,y_2-y_1)=0$. En el espacio la ecuación implícita de una sola ecuación describe un plano.

\begin{definicion}[Plano]\label{def:plano}
El \emph{plano} que pasa por $P_0=(x_0,y_0,z_0)$ con \emph{vector normal} $\vn=(a,b,c)\neq\vcero$ es el conjunto de los puntos $X=(x,y,z)$ tales que $\vn\cdot(X-P_0)=0$, es decir,
\[
a(x-x_0)+b(y-y_0)+c(z-z_0)=0\quad\Longleftrightarrow\quad ax+by+cz=d,\qquad d=\vn\cdot P_0 .
\]
\end{definicion}

Recíprocamente, toda ecuación $ax+by+cz=d$ con $(a,b,c)\neq\vcero$ es un plano de normal $(a,b,c)$: basta tomar uno de sus puntos como $P_0$. Para el plano que pasa por tres puntos no alineados $P$, $Q$, $R$ se toma como normal $\vn=\overrightarrow{PQ}\times\overrightarrow{PR}$, que es perpendicular a ambos vectores (\cref{prop:vectorial}).

\begin{ejemplo}\label{ej:t2-plano}
El plano que pasa por $(1,0,0)$, $(0,2,0)$ y $(0,0,3)$ tiene normal $(-1,2,0)\times(-1,0,3)=(6,3,2)$ y ecuación $6x+3y+2z=6$, es decir, $\frac x1+\frac y2+\frac z3=1$: los denominadores son los puntos de corte con los ejes.
\end{ejemplo}

\begin{proposicion}[Posición relativa de dos planos]\label{prop:dos-planos}
Sean $\Pi_1:\vn_1\cdot X=d_1$ y $\Pi_2:\vn_2\cdot X=d_2$.
\begin{enumerate}
  \item Si $\vn_1$ y $\vn_2$ son proporcionales, los planos son paralelos: coinciden si las ecuaciones son proporcionales y no tienen puntos comunes en caso contrario.
  \item Si no lo son, se cortan en una recta de vector director $\vn_1\times\vn_2$.
\end{enumerate}
El \emph{ángulo} entre $\Pi_1$ y $\Pi_2$ es el ángulo agudo (o recto) $\varphi$ entre sus normales: $\cos\varphi=\abs{\vn_1\cdot\vn_2}/(\norm{\vn_1}\norm{\vn_2})$.
\end{proposicion}

\begin{proof}[Idea de la demostración]
(i) Si $\vn_2=\lambda\vn_1$, todo punto común cumpliría $d_2=\vn_2\cdot X=\lambda d_1$; si $d_2=\lambda d_1$ las ecuaciones son equivalentes, y si no, son incompatibles. (ii) Una dirección $\vv$ contenida en ambos planos es perpendicular a $\vn_1$ y a $\vn_2$; $\vn_1\times\vn_2\neq\vcero$ tiene esa propiedad, y la existencia de un punto común se obtiene resolviendo el sistema (\cref{sec:t2-gauss}).
\end{proof}

\begin{ejemplo}\label{ej:t2-recta-interseccion}
Los planos $x+y+z=1$ y $x-y=0$ tienen normales $(1,1,1)$ y $(1,-1,0)$, no proporcionales. La recta común tiene dirección $(1,1,1)\times(1,-1,0)=(1,1,-2)$; un punto común es $(0,0,1)$ (con $x=y=0$). Luego la recta es $\vr(t)=(t,\,t,\,1-2t)$; en efecto, $t+t+1-2t=1$ y $t-t=0$.
\end{ejemplo}

Para cortar una recta $\vr(t)=P_0+t\vv$ con un plano $\vn\cdot X=d$ se sustituye $\vr(t)$ en la ecuación del plano y se despeja $t$. Por ejemplo, la recta $(1+t,\,t,\,t)$ y el plano $x+2y+3z=7$ dan $1+6t=7$, $t=1$, y el punto de corte es $(2,1,1)$. Si $\vn\cdot\vv=0$, la recta es paralela al plano: o no lo corta o está contenida en él.

\section{Sistemas lineales y eliminación gaussiana}\label{sec:t2-gauss}

\begin{definicion}[Sistema lineal y matriz]\label{def:sistema}
Un sistema de $m$ ecuaciones lineales con $n$ incógnitas ($m,n\le3$ en este tema) se escribe $A\vx=\vb$, donde $A=(a_{ij})$ es la \emph{matriz de coeficientes}, $\vx$ el vector de incógnitas y $\vb$ el de términos independientes. El producto $A\vx$ es el vector cuya componente $i$ es el producto escalar de la fila $i$ de $A$ por $\vx$. El sistema es \emph{compatible determinado} si tiene una única solución, \emph{compatible indeterminado} si tiene infinitas e \emph{incompatible} si no tiene ninguna.
\end{definicion}

Más en general, el producto de matrices $AB$ (definido cuando el número de columnas de $A$ coincide con el de filas de $B$) tiene como elemento $(i,j)$ el producto escalar de la fila $i$ de $A$ por la columna $j$ de $B$. La matriz identidad $I$ tiene unos en la diagonal y ceros fuera, y cumple $AI=IA=A$.

\begin{definicion}[Operaciones elementales]\label{def:elementales}
Las \emph{operaciones elementales} sobre las ecuaciones de un sistema son: (E1) intercambiar dos ecuaciones; (E2) multiplicar una ecuación por un número no nulo; (E3) sumar a una ecuación un múltiplo de otra.
\end{definicion}

\begin{proposicion}\label{prop:elementales}
Las operaciones elementales no cambian el conjunto de soluciones de un sistema.
\end{proposicion}

\begin{proof}
Basta verlo para (E3), pues (E1) y (E2) son evidentes (para (E2), porque se puede deshacer dividiendo por el mismo número no nulo). Sean $L_1(\vx)=b_1$ y $L_2(\vx)=b_2$ dos ecuaciones y sustitúyase la segunda por $L_2(\vx)+\lambda L_1(\vx)=b_2+\lambda b_1$. Toda solución del sistema original cumple la nueva ecuación, sumando $\lambda$ veces la primera a la segunda. Recíprocamente, la operación se deshace restando $\lambda$ veces la primera ecuación, de modo que toda solución del sistema nuevo lo es del original.
\end{proof}

El \emph{método de Gauss} aplica operaciones elementales hasta llegar a un sistema \emph{escalonado} (cada ecuación tiene al menos una incógnita inicial menos que la anterior) y lo resuelve por \emph{sustitución regresiva}, de abajo arriba. En la práctica se trabaja con la \emph{matriz ampliada} $(A\mid\vb)$.

\begin{ejemplo}\label{ej:t2-gauss}
Se resuelve $x+2y+z=4$, $2x+5y+3z=10$, $-x+y+4z=4$:
\[
\left(\begin{array}{rrr|r}1&2&1&4\\2&5&3&10\\-1&1&4&4\end{array}\right)
\xrightarrow[F_3+F_1]{F_2-2F_1}
\left(\begin{array}{rrr|r}1&2&1&4\\0&1&1&2\\0&3&5&8\end{array}\right)
\xrightarrow{F_3-3F_2}
\left(\begin{array}{rrr|r}1&2&1&4\\0&1&1&2\\0&0&2&2\end{array}\right).
\]
De abajo arriba: $z=1$, $y=2-z=1$, $x=4-2y-z=1$. La solución es única: $(1,1,1)$.
\end{ejemplo}

Si al escalonar aparece una fila $(0\ 0\ 0\mid c)$ con $c\neq0$, el sistema es incompatible; si aparece una fila entera de ceros y quedan menos ecuaciones que incógnitas, las incógnitas que no encabezan ninguna fila son \emph{libres} y el sistema es compatible indeterminado. Geométricamente, en un sistema $3\times3$ cada ecuación es un plano: la solución única corresponde a tres planos con un único punto común; la indeterminación, a planos que comparten una recta (o un plano); la incompatibilidad, a planos sin punto común a los tres.

\section{Determinantes de orden 2 y 3}

\begin{definicion}[Determinantes]\label{def:det3}
\[
\det\begin{pmatrix}a&b\\c&d\end{pmatrix}=\begin{vmatrix}a&b\\c&d\end{vmatrix}=ad-bc,
\qquad
\begin{vmatrix}a_1&a_2&a_3\\b_1&b_2&b_3\\c_1&c_2&c_3\end{vmatrix}
=a_1\begin{vmatrix}b_2&b_3\\c_2&c_3\end{vmatrix}-a_2\begin{vmatrix}b_1&b_3\\c_1&c_3\end{vmatrix}+a_3\begin{vmatrix}b_1&b_2\\c_1&c_2\end{vmatrix}.
\]
\end{definicion}

Comparando con la \cref{def:vectorial}, el determinante $3\times3$ de filas $\va,\vb,\mathbf c$ es exactamente el producto mixto $\va\cdot(\vb\times\mathbf c)$.

\begin{teorema}[Interpretación geométrica]\label{teo:det-area}
\begin{enumerate}
  \item $\abs{ad-bc}$ es el área del paralelogramo de lados $(a,b)$ y $(c,d)$ en $\R^2$.
  \item El valor absoluto del determinante $3\times3$ de filas $\va,\vb,\mathbf c$ es el volumen del paralelepípedo de aristas $\va,\vb,\mathbf c$.
\end{enumerate}
\end{teorema}

\begin{proof}
(i) Se ven los vectores en $\R^3$ como $(a,b,0)$ y $(c,d,0)$. Su producto vectorial es $(b\cdot0-0\cdot d,\ 0\cdot c-a\cdot0,\ ad-bc)=(0,0,ad-bc)$, cuya norma, por la \cref{prop:vectorial}(iv), es el área del paralelogramo. (ii) Es la \cref{prop:volumen}, dado que el determinante es el producto mixto.
\end{proof}

\begin{proposicion}[Propiedades de los determinantes]\label{prop:det-propiedades}
Para matrices $2\times2$ y $3\times3$:
\begin{enumerate}
  \item intercambiar dos filas cambia el signo del determinante; si dos filas son iguales, el determinante es nulo;
  \item multiplicar una fila por $\lambda$ multiplica el determinante por $\lambda$;
  \item sumar a una fila un múltiplo de otra no cambia el determinante;
  \item $\det(A^{\mathsf T})=\det A$ y $\det(AB)=\det A\cdot\det B$.
\end{enumerate}
\end{proposicion}

Para orden~2 todas las propiedades se comprueban con la fórmula (véase el \cref{eje:2-producto-det}); para orden~3 se deducen de las propiedades del producto mixto, y la demostración general, para cualquier orden, se da en el módulo~3 del curso. La propiedad~(iii) dice que la eliminación gaussiana con operaciones (E3) conserva el determinante: el de una matriz escalonada es el producto de los elementos de su diagonal. En el \cref{ej:t2-gauss}, $\det A=1\cdot1\cdot2=2$.

\begin{teorema}[Sistemas $2\times2$ y regla de Cramer]\label{teo:cramer}
El sistema $ax+by=e$, $cx+dy=f$ tiene solución única si y solo si $\Delta=ad-bc\neq0$, y en ese caso
\[
x=\frac{1}{\Delta}\begin{vmatrix}e&b\\f&d\end{vmatrix},\qquad y=\frac{1}{\Delta}\begin{vmatrix}a&e\\c&f\end{vmatrix}.
\]
\end{teorema}

\begin{proof}
Multiplicando la primera ecuación por $d$, la segunda por $b$ y restando, se elimina $y$: $(ad-bc)x=ed-bf$. Análogamente, multiplicando por $c$ y por $a$ y restando, $(ad-bc)y=af-ce$. Si $\Delta\neq0$ estas igualdades determinan $x$ e $y$, y se comprueba por sustitución que son solución. Si $\Delta=0$, las filas $(a,b)$ y $(c,d)$ son proporcionales (el paralelogramo que forman tiene área nula, \cref{teo:det-area}), las dos rectas son paralelas o coincidentes y no hay solución única.
\end{proof}

Lo mismo vale en orden~3: un sistema $3\times3$ tiene solución única si y solo si $\det A\neq0$, es decir, si las tres normales (las filas de $A$) no son coplanarias (\cref{prop:volumen}); la regla de Cramer sustituye la columna $j$ de $A$ por $\vb$ para obtener la incógnita $j$.

\begin{definicion}[Matriz inversa]\label{def:inversa}
Una matriz cuadrada $A$ es \emph{invertible} si existe $A^{-1}$ con $AA^{-1}=A^{-1}A=I$. Entonces el sistema $A\vx=\vb$ tiene la única solución $\vx=A^{-1}\vb$.
\end{definicion}

Para orden~2, si $ad-bc\neq0$,
\begin{equation}\label{eq:t2-inversa}
\begin{pmatrix}a&b\\c&d\end{pmatrix}^{-1}=\frac{1}{ad-bc}\begin{pmatrix}d&-b\\-c&a\end{pmatrix},
\end{equation}
como se comprueba multiplicando. Una matriz es invertible si y solo si su determinante es no nulo.

\begin{observacion}[Errores frecuentes]
\begin{enumerate}[label=(\alph*)]
\item La ecuación $ax+by+cz=d$ en el espacio es un plano, no una recta; una recta del espacio necesita dos ecuaciones implícitas o una paramétrica.
\item En la eliminación, cada fila nueva debe calcularse a partir de las filas del paso \emph{anterior}; mezclar filas ya modificadas en el mismo paso es una fuente habitual de errores.
\item $\det(A+B)\neq\det A+\det B$ y $\det(\lambda A)=\lambda^n\det A$ en orden $n$, no $\lambda\det A$.
\item Un determinante nulo no significa «sin solución»: el sistema puede ser incompatible o tener infinitas soluciones.
\end{enumerate}
\end{observacion}

\begin{observacion}[Conexión con el curso]
Los sistemas lineales vuelven en el módulo~3 (inversa, determinante y rango, tema~2) y, sobre todo, en la regresión lineal del módulo~8, cuyo ajuste por mínimos cuadrados se reduce a resolver las \emph{ecuaciones normales} $X^{\mathsf T}X\mathbf w=X^{\mathsf T}\mathbf y$. El determinante como factor de cambio de área reaparece en el cambio de variables de las integrales dobles (\cref{tema:6}), donde el elemento de área en polares es $r\,dr\,d\theta$.
\end{observacion}

\begin{lecturas}
\item MIT 18.02SC~\cite{mit1802}, unidad~1, parte~A, sesiones~5, 6 y~8 (determinantes, volúmenes, planos), parte~B, sesiones~9--14 (matrices y sistemas) y parte~C, sesiones~15--16 (rectas e intersecciones).
\item OpenStax \emph{Calculus}, vol.~3~\cite{openstax3}, §2.5 «Equations of Lines and Planes in Space».
\item Strang, \emph{Calculus}~\cite{strangcalc}, §11.2 «Planes and Projections» y §11.4 «Matrices and Linear Equations»; Strang, \emph{Introduction to Linear Algebra}~\cite{strangla}, §2.2.
\end{lecturas}

\section*{Ejercicios del tema 2}
\addcontentsline{toc}{section}{Ejercicios del tema 2}

\begin{ejercicio}{1}{eje:2-posicion}
Sin resolver, indique la posición relativa de cada par de planos y, en consecuencia, si el sistema formado por sus dos ecuaciones es compatible determinado, indeterminado o incompatible: (a) $x+2y-z=1$ y $2x+4y-2z=5$; (b) $x+2y-z=1$ y $3x+6y-3z=3$; (c) $x+y=1$ e $y+z=1$.
\end{ejercicio}

\begin{ejercicio}{2}{eje:2-planos}
Halle la ecuación del plano: (a) que pasa por $(2,0,-1)$ y es perpendicular a $\vi+2\vj-2\vk$; (b) que pasa por $(1,0,1)$, $(2,-1,2)$ y $(-1,3,2)$.
\fuente{adaptado de MIT 18.02SC~\cite{mit1802}, problemas suplementarios, 1E-1.}
\end{ejercicio}

\begin{ejercicio}{2}{eje:2-diedro}
Calcule el ángulo que forman los planos $2x-y+z=3$ y $x+y+2z=1$.
\fuente{adaptado de MIT 18.02SC~\cite{mit1802}, problemas suplementarios, 1E-2.}
\end{ejercicio}

\begin{ejercicio}{2}{eje:2-interseccion}
(a) Escriba la recta que pasa por $(2,-1,-1)$ y es perpendicular al plano $x-y+2z=3$. (b) Halle el punto en que la recta que pasa por $(0,1,2)$ y $(2,0,3)$ corta al plano $x+4y+z=4$.
\fuente{adaptado de MIT 18.02SC~\cite{mit1802}, problemas suplementarios, 1E-3 y 1E-4.}
\end{ejercicio}

\begin{ejercicio}{2}{eje:2-gauss}
Resuelva por eliminación gaussiana el sistema $x_1-x_2+x_3=2$, $x_2+x_3=0$, $-x_1-x_2+2x_3=3$, y calcule el determinante de su matriz de coeficientes.
\fuente{adaptado de MIT 18.02SC~\cite{mit1802}, problemas suplementarios, 1G-3.}
\end{ejercicio}

\begin{ejercicio}{2}{eje:2-gauss-casos}
Resuelva y clasifique: (a) $x+y+z=6$, $2x-y+z=3$, $x+2y-z=2$; (b) $x+y+z=1$, $x+2y+3z=2$, $2x+3y+4z=3$; (c) el sistema (b) con el último término independiente cambiado por~4. Interprete geométricamente cada caso.
\end{ejercicio}

\begin{ejercicio}{2}{eje:2-parametro}
(a) ¿Para qué valor de $c$ tiene el sistema homogéneo $x_1-x_2+x_3=0$, $2x_1+x_2+x_3=0$, $-x_1+cx_2+2x_3=0$ soluciones distintas de la trivial? (b) Para ese valor, halle una solución no trivial usando el producto vectorial de las dos primeras filas y justifique por qué funciona.
\fuente{adaptado de MIT 18.02SC~\cite{mit1802}, problemas suplementarios, 1H-3.}
\end{ejercicio}

\begin{ejercicio}{2}{eje:2-areas}
Calcule con determinantes el área del triángulo de vértices (a) $(0,0)$, $(1,2)$, $(1,-1)$; (b) $(1,2)$, $(1,-1)$, $(2,3)$.
\fuente{adaptado de MIT 18.02SC~\cite{mit1802}, problemas suplementarios, 1C-3.}
\end{ejercicio}

\begin{ejercicio}{2}{eje:2-inversa}
Sea $A=\begin{pmatrix}4&3\\3&2\end{pmatrix}$. Calcule $A^{-1}$ con la fórmula~\eqref{eq:t2-inversa}, compruebe que $AA^{-1}=I$ y resuelva $A\vx=(2,3)$.
\fuente{adaptado de MIT 18.02SC~\cite{mit1802}, problemas suplementarios, 1G-2.}
\end{ejercicio}

\begin{ejercicio}{2}{eje:2-produccion}
Con tres materias primas $M_1,M_2,M_3$ se fabrican tres productos. Cada unidad de $P_1$ consume $1$, $2$ y $3$ unidades de $M_1$, $M_2$, $M_3$; cada unidad de $P_2$, $1$, $3$ y $5$; cada unidad de $P_3$, $3$, $5$ y $8$. En una jornada se consumieron $137$, $279$ y $448$ unidades de $M_1$, $M_2$, $M_3$. (a) Plantee el sistema $A\vx=\vb$ y calcule $\det A$. (b) Halle cuántas unidades de cada producto se fabricaron. (c) Si $P_3$ consumiera $2$, $5$ y $8$ unidades, ¿cuánto valdría $\det A$ y qué ocurriría con los datos de consumo dados?
\fuente{adaptado de MIT 18.02SC~\cite{mit1802}, \emph{problem set}~2, parte~II, problema~2.}
\end{ejercicio}

\begin{ejercicio}{2}{eje:2-cramer}
Resuelva por la regla de Cramer el sistema $3x-y+z=1$, $-x+2y+z=2$, $x-y+z=-3$, y compruebe el resultado por sustitución.
\fuente{adaptado de MIT 18.02SC~\cite{mit1802}, problemas suplementarios, 1H-1.}
\end{ejercicio}

\begin{ejercicio}{3}{eje:2-distancia}
(a) Demuestre que la distancia del punto $P_1=(x_1,y_1,z_1)$ al plano $ax+by+cz=d$ es $\dfrac{\abs{ax_1+by_1+cz_1-d}}{\sqrt{a^2+b^2+c^2}}$. (Indicación: proyecte $\overrightarrow{P_0P_1}$, con $P_0$ en el plano, sobre la normal.) (b) Deduzca la distancia del origen al plano y calcule la distancia de $(1,2,3)$ al plano $2x-y+2z=10$.
\fuente{adaptado de MIT 18.02SC~\cite{mit1802}, problemas suplementarios, 1E-6.}
\end{ejercicio}

\begin{ejercicio}{3}{eje:2-vandermonde}
Demuestre que
$\begin{vmatrix}1&1&1\\x_1&x_2&x_3\\x_1^2&x_2^2&x_3^2\end{vmatrix}=(x_1-x_2)(x_2-x_3)(x_3-x_1)$
y deduzca que, si $x_1,x_2,x_3$ son distintos, existe una única parábola $y=c_0+c_1x+c_2x^2$ que pasa por tres puntos dados $(x_i,y_i)$.
\fuente{adaptado de MIT 18.02SC~\cite{mit1802}, problemas suplementarios, 1C-4.}
\end{ejercicio}

\begin{ejercicio}{3}{eje:2-compatibilidad}
Sea el sistema $a_1x+b_1y=c_1$, $a_2x+b_2y=c_2$ con $a_1b_2-a_2b_1=0$ y $a_1\neq0$. Demuestre que es compatible si y solo si $c_2=\frac{a_2}{a_1}c_1$, y que en ese caso tiene infinitas soluciones.
\fuente{adaptado de MIT 18.02SC~\cite{mit1802}, problemas suplementarios, 1H-5.}
\end{ejercicio}

\begin{ejercicio}{3}{eje:2-producto-det}
Sean $A$ y $B$ matrices $2\times2$. (a) Demuestre que $\det(AB)=\det A\cdot\det B$. (b) Demuestre, a partir de la definición de inversa, que si $A$ y $B$ son invertibles, entonces $AB$ lo es y $(AB)^{-1}=B^{-1}A^{-1}$. (c) Deduzca que $\det(A^{-1})=1/\det A$.
\fuente{adaptado de MIT 18.02SC~\cite{mit1802}, problemas suplementarios, 1G-5.}
\end{ejercicio}
